\documentclass[11pt, a4paper]{article}

\usepackage[utf8]{inputenc}
\usepackage[T1]{fontenc}
\usepackage[english]{babel}
\usepackage{amsmath, amssymb, amsthm}
\usepackage[margin=1.1in]{geometry}
\usepackage{hyperref}
\usepackage{xcolor}
\usepackage{listings}
\usepackage{booktabs}

\hypersetup{colorlinks=true, linkcolor=blue!70!black, citecolor=red!70!black}

\definecolor{backcolor}{rgb}{0.95, 0.95, 0.92}
\lstset{
    basicstyle=\ttfamily\footnotesize,
    breaklines=true,
    frame=single,
    backgroundcolor=\color{backcolor},
    keywordstyle=\bfseries\color{blue!70!black}
}

\title{\textbf{On the Physical Vacuity of Manufactured Singularities:}\\
\Large An Epistemic Audit of the OpenAI Navier-Stokes Formalization}
\author{
    \textbf{The MechanicaFluidorum Program} \\
    \textit{Socrate AI Lab} \\
    \small Non-Profit Research Organization (French Association Loi 1901 for Neuro-Symbolic Scientific AI) \\
    \small \href{https://github.com/xaviercallens/OpenAI-NSE-Epistemic-Audit}{\texttt{github.com/xaviercallens/OpenAI-NSE-Epistemic-Audit}} \\
    \small Certified Archive: \href{https://doi.org/10.5281/zenodo.22696718}{DOI: 10.5281/zenodo.22696718}
}
\date{September 2026}

\begin{document}

\maketitle

\begin{abstract}
In September 2026, an OpenAI multi-agent system formalized finite-time blow-up proofs for the forced 3D Navier-Stokes equations (Millennium Prize Alternatives C and D) and the unforced 3D Euler equations within the Lean 4 proof assistant. While this represents a landmark achievement in automated theorem proving, a rigorous epistemic audit of the underlying codebase reveals that the singularities are driven by pathological, manufactured mathematics rather than natural autonomous fluid dynamics. For the forced Navier-Stokes case, the external force is reverse-engineered as the exact residual of a pre-constructed singular trajectory, artificially canceling viscous thermodynamic dissipation. For the unforced Euler equations, the initial condition relies on an infinite fractal induction of vortex packets pushed into unbounded ultraviolet frequencies ($\kappa_n \to \infty$) at $t=0$, violating the fundamental continuum hypothesis of fluid mechanics well below the Kolmogorov microscale. We demonstrate that while syntax-perfect within abstract Sobolev spaces, these proofs exploit the infinite divisibility of the classical continuum and leave the physical Millennium Problem entirely open.
\end{abstract}

\section{Introduction: Syntax vs. Semantics in Formal Proofs}
The application of large-scale AI agents to formalize mathematics in Lean 4 guarantees absolute logical consistency. The repository released by OpenAI \cite{OpenAI2026} (audited at commit hash \texttt{8937a8f4cbc7abaab5e9e97d1cc7f5d2319d9538}) verifies that under specific engineered conditions, fluid equations can develop singularities. However, in mathematical physics, a formal proof is only as physically relevant as the phenomenological realism of its initial conditions and forcing functions. As articulated by Batchelor \cite{Batchelor1967} and Landau \& Lifshitz \cite{Landau1987}, fluid dynamics is subordinate to the physical hypotheses that constrain its PDEs. In this brief communication, we reverse-engineer the Lean 4 codebase to expose the ``manufactured'' nature of these singularities.

\section{The Forced Navier-Stokes Singularity: A Tautological Syringe}
To prove blow-up with viscosity $\nu > 0$, the proof targets the forced scenarios (Alternatives C and D). An inspection of \texttt{NavierStokes/ProblemStatement.lean} reveals the fluid starts from absolute rest (\texttt{zero\_initial\_velocity}). The singularity is entirely driven by the external force $f(x,t)$.

Crucially, in \texttt{NavierStokes/R3/ProblemStatement.lean} and \texttt{CandidateFromLimits.lean}, the forcing $f(x,t)$ is explicitly defined as the exact residual of a singularity-bound activated field $\mathbf{u}_{\text{act}}$:
\begin{equation}
    f(x,t) = \partial_t \mathbf{u}_{\text{act}} + (\mathbf{u}_{\text{act}} \cdot \nabla)\mathbf{u}_{\text{act}} - \nu \Delta \mathbf{u}_{\text{act}} + \nabla p_{\text{act}}
\end{equation}

In the PDE literature, this is known as the \textit{Method of Manufactured Solutions} (MMS) \cite{Roache2002}. It is a tool for code verification, not physical discovery. Viscous dissipation ($-\nu\Delta \mathbf{u}$) is the thermodynamic mechanism by which kinetic energy is irreversibly transferred into thermal entropy at small scales \cite{Kolmogorov1941}. By defining the force as the residual of a collapsing vortex, the framework acts as a non-autonomous ``Maxwell's Demon,'' injecting targeted momentum with infinite precision to artificially cancel the viscous dissipation. 

It must be explicitly acknowledged that Charles Fefferman's official Clay Mathematics Institute formulation \cite{Fefferman2000} for Alternatives C and D explicitly \textit{permits} any globally defined, smooth external force $f(x,t)$. The AI did not cheat; it perfectly optimized for the highly permissive, purely mathematical bounds set by the CMI. The epistemological failure lies just as much in the pure-math wording of the Millennium Prize itself as it does in the AI's methodology. The AI exposed a misalignment between the mathematical framing and the physical reality. While the OpenAI formalization flawlessly satisfies the highly permissive mathematical wording of Millennium Prize Statements C and D, the singularity is deterministically forced by an artificial construct. This yields no insight into autonomous fluid stability and implies that an omniscient forcing is required to overcome natural viscous dissipation.

\section{The Unforced Euler Singularity: The Fractal Ultraviolet Bomb}
For the unforced Euler equations, the singularity must emerge solely from the initial condition $\mathbf{u}_0(x)$. The audit of \texttt{Euler/PacketInitialSmoothLimit.lean} reveals $\mathbf{u}_0$ is constructed as the limit of an infinite series of vortex packet increments evaluated at spatial frequencies $\kappa_n$:
\begin{equation}
    \mathbf{u}_0(\mathbf{x}) = \sum_{n=1}^\infty \Big(\mathbf{v}_{n,\text{high}}(\kappa_n, \mathbf{x}) + \mathbf{v}_{n,\text{mean}}(\kappa_n, \mathbf{x})\Big)
\end{equation}

As defined in \texttt{Euler/PacketSourceScaleChoice.lean} and \texttt{PacketSourceScaleSequence.lean}, the scale sequence grows factorially ($s_{n+1} \propto (J+n)^2 s_n$), meaning $\kappa_n \to \infty$ explodes along a factorial-exponential tower. Injecting active kinetic energy into spatial scales approaching zero at $t=0$ implies the existence of coherent vortex structures smaller than atomic or even Planck scales. The Navier-Stokes and Euler equations are macroscopic continuum idealizations valid \textit{if and only if} the Knudsen number is small ($Kn = \lambda_{\text{mfp}}/L \ll 1$) \cite{Batchelor1967, Landau1987}. 

To satisfy the Fefferman criteria \cite{Fefferman2000}, which mandate that the initial condition $\mathbf{u}_0(x)$ must be infinitely differentiable ($C^\infty$) with bounded kinetic energy, the AI relies on a carefully tuned amplitude-frequency trade-off. The amplitude of the vortex packets is scaled inversely to their frequency (i.e., the energy spectrum decays rapidly), ensuring that the infinite sum of the norms converges to a finite value at $t=0$. Mathematically, this infinite induction successfully converges within abstract Sobolev spaces $H^m(\mathbb{R}^3)$, satisfying the rigorous bounds of the Lean 4 kernel. However, epistemologically, the Navier-Stokes and Euler PDEs are macroscopic asymptotic models. By operating strictly outside these bounds, this initial condition violates the physical applicability limit of the equations.

\section{Conclusion}
The OpenAI formalization definitively closes the specific wording of Alternatives C and D of the Millennium Prize \cite{Fefferman2000}, proving that the mathematical continuum ($\mathbb{R}^3$) contains enough infinite divisibility to allow a reverse-engineered force to break the equations. However, from a phenomenological standpoint, these artifacts highlight why classical continuum mechanics breaks down at ultraviolet scales. The global regularity of natural, unforced 3D Navier-Stokes equations (Statements A and B) remains profoundly open, highlighting the necessity for mechanisms of geometric or topological regularization at the ultraviolet scale.

\begin{thebibliography}{9}

\bibitem{OpenAI2026}
OpenAI. \textit{Finite time blowup for Navier-Stokes and Euler equations}. Lean 4 Formalization Repository (\url{https://github.com/openai/NavierStokesAndEuler}), audited at commit \texttt{8937a8f4cbc7abaab5e9e97d1cc7f5d2319d9538}, 2026.

\bibitem{Batchelor1967}
G. K. Batchelor. \textit{An Introduction to Fluid Dynamics}. Cambridge University Press (1967).

\bibitem{Landau1987}
L. D. Landau, \& E. M. Lifshitz. \textit{Fluid Mechanics}, Course of Theoretical Physics, Vol. 6. Pergamon Press (1987).

\bibitem{Kolmogorov1941}
A. N. Kolmogorov. \textit{The local structure of turbulence in incompressible viscous fluid for very large Reynolds numbers}. Dokl. Akad. Nauk SSSR, 30, 301-305 (1941).

\bibitem{Roache2002}
P. J. Roache. \textit{Code Verification by the Method of Manufactured Solutions}. Journal of Fluids Engineering, 124(1), 4-10 (2002).

\bibitem{Fefferman2000}
C. Fefferman. \textit{Existence and smoothness of the Navier-Stokes equation}. Clay Mathematics Institute Millennium Prize Problems (2000).

\end{thebibliography}

\end{document}
